\section{V. ĐỀ XUẤT KHUYẾN NGHỊ}

\subsection{5.1. Đối với các cơ sở giáo dục đại học}
Các cơ sở giáo dục đại học cần ưu tiên đầu tư chuẩn hóa hạ tầng dữ liệu, nâng cấp Hệ thống Quản lý Học tập để hỗ trợ ghi vết dữ liệu sự kiện có độ chi tiết cao theo chuẩn quốc tế xAPI hoặc Caliper Analytics.

Đồng thời, cần hướng tới thành lập Trung tâm Phân tích Học tập để tích hợp mô hình đánh giá quy trình vào bảng điều khiển (Dashboard) dành cho cố vấn học tập và giảng viên, giúp chủ động nhận diện sinh viên cần hỗ trợ.

\subsection{5.2. Đối với giảng viên và người thiết kế khóa học}
Giảng viên cần chủ động xây dựng các Mô hình quy trình tham chiếu thông qua việc thiết kế kịch bản học tập rõ ràng, chuẩn hóa thứ tự các hoạt động và mốc tiến trình để tạo cơ sở đối sánh chính xác cho mô hình Process Mining.

\subsection{5.3. Đối với các đơn vị phát triển giải pháp EdTech}
Các đơn vị phát triển phần mềm giáo dục cần đẩy mạnh tích hợp Module Process Mining trực tiếp vào các hệ thống LMS thông qua việc phát triển các Plugin hoặc API kết nối với Moodle hay Canvas, giúp tự động hóa việc đưa ra gợi ý lộ trình học tập cá nhân hóa cho người học.
